\paragrafos{Διαίρεση ρητών: μοιράζω και τα χρέη}{Α.7.6, «Διαίρεση ρητών αριθμών»}

\begin{anagke}
Στο τέλος της εκδρομής ο Νίκος, που κρατάει το ταμείο, ανακοινώνει με σκυμμένο κεφάλι: «Το ταμείο έμεινε στο $-24$ ευρώ. Χρωστάμε $24$ ευρώ». Η Μυρτώ ρωτάει: «Αν το μοιράσουμε εξίσου στα έξι παιδιά της παρέας, πόσο ήταν το καθένα;». Η Πολυξένη το γράφει: $(-24):6$. Ο Βασίλης λέει «τέσσερα» και ο Στέλιος διαμαρτύρεται: «Πόσο, χρέος ή κέρδος;». Αφού το χρέος μοιράζεται, κάθε παιδί χρωστά $4$ ευρώ, δηλαδή το πηλίκο είναι $-4$.

Ο Ηλίας ρωτάει το αντίστροφο: «Αν κάθε παιδί πρέπει να βάλει $3$ ευρώ ($-3$ στο ταμείο του) και το σύνολο είναι $-24$, πόσα παιδιά είναι;». Ο Νίκος σκέφτεται: «Πόσες φορές χωράει το $-3$ στο $-24$; Οκτώ φορές». Άρα $(-24):(-3)=+8$. Η Αγγελική τα βάζει σε εξίσωση: $(-3)\cdot x=-24$ και $x=8$. Και η Πολυξένη θυμάται ότι η διαίρεση είναι ο αντίστροφος δρόμος του πολλαπλασιασμού, όπως η αφαίρεση της πρόσθεσης.

Από τον πολλαπλασιασμό ήδη ξέρουν τους κανόνες των προσήμων. Τους ίδιους κανόνες θα βρουν και για τη διαίρεση, και θα δουν γιατί, πάλι, δεν μπορούν να διαιρέσουν με το $0$.
\end{anagke}

\begin{lexi}{διαίρεση, πηλίκο, αντίστροφος}{διαίρεση, πηλίκο, λόγος}
Η \lexiword{διαίρεση} είναι το «δι-αιρώ», δηλαδή «χωρίζω σε μέρη». Ο \lexiword{διαιρετέος} είναι αυτός που μοιράζεται, ο \lexiword{διαιρέτης} αυτός που μοιράζει και το \lexiword{πηλίκο} (από το «πηλίκος», «πόσος») το αποτέλεσμα. Ο \lexiword{λόγος} δύο ρητών είναι το πηλίκο τους, γραμμένο ως κλάσμα.
\end{lexi}

\begin{metafrasi}
Η διαίρεση $\alpha:\beta$ είναι η λύση της εξίσωσης. & $\beta\cdot x=\alpha$ (με $\beta\neq0$)\\
Μοιράζω χρέος $24$ € σε $6$ παιδιά. & $(-24):6=-4$\\
Πόσες φορές χωράει το $-3$ στο $-24$; & $(-24):(-3)=+8$\\
Διαίρεση σημαίνει πολλαπλασιασμός με τον αντίστροφο. & $\alpha:\beta=\alpha\cdot\dfrac{1}{\beta}$\\
Κλάσμα με αρνητικό. & $\dfrac{-6}{3}=\dfrac{6}{-3}=-2$\\
\end{metafrasi}

\tmima{Η θεωρία}

\begin{orismos}{πηλίκο}
\textbf{Πηλίκο} $\alpha:\beta$ δύο ρητών με $\beta\neq0$ ορίζεται ως η μοναδική λύση της εξίσωσης $\beta\cdot x=\alpha$.
\end{orismos}

Για παράδειγμα το $12:(-3)$ είναι η λύση της $(-3)\cdot x=12$, δηλαδή $x=-4$. Από αυτό προκύπτει ότι τα πρόσημα είναι τα ίδια με του πολλαπλασιασμού:

\begin{kanonas}{διαίρεση ρητών}
Το πηλίκο δύο ρητών έχει απόλυτη τιμή το πηλίκο των απόλυτων τιμών και \textbf{πρόσημο} θετικό αν είναι ομόσημοι και αρνητικό αν είναι ετερόσημοι. Δηλαδή $\alpha:\beta$ και $\alpha\cdot\beta$ είναι \textbf{ομόσημοι}. Το $0:\beta=0$. Δεν διαιρούμε με το $0$.
\end{kanonas}

Έτσι $(-24):6=-4$, $(-24):(-3)=+8$, $18:(-6)=-3$, $0:(-5)=0$.

\tmimasmall{Διαίρεση και αντίστροφος}

Ο \textbf{αντίστροφος} ενός ρητού $\alpha\neq0$ είναι ο $\frac{1}{\alpha}$ (γινόμενο $1$). Ο αντίστροφος του $-4$ είναι ο $-\frac{1}{4}$, του $-\frac{2}{3}$ ο $-\frac{3}{2}$. Το $0$ δεν έχει αντίστροφο. Κάθε διαίρεση γίνεται πολλαπλασιασμός: $\alpha:\beta=\alpha\cdot\frac{1}{\beta}$, όπως και στα κλάσματα. Έτσι $\left(-\frac{3}{4}\right):\frac{1}{2}=\left(-\frac{3}{4}\right)\cdot2=-\frac{3}{2}$.

\tmimasmall{Το πηλίκο ως λόγος}

Το πηλίκο δύο ρητών γράφεται ως κλάσμα (\emph{λόγος}): $\dfrac{-6}{3}=-2$, $\dfrac{6}{-3}=-2$, $\dfrac{-6}{-3}=+2$. Ισχύει $\dfrac{-\alpha}{\beta}=\dfrac{\alpha}{-\beta}=-\dfrac{\alpha}{\beta}$. Τα σύνθετα κλάσματα λύνονται όπως πριν: $\dfrac{-\frac{6}{5}}{\frac{3}{10}}=-\dfrac{6}{5}\cdot\dfrac{10}{3}=-4$.

\begin{prosoxi}
\begin{itemize}[leftmargin=1.2em,itemsep=0.3ex,topsep=2pt]
\item Δεν υπάρχει $\alpha:0$: η $0\cdot x=\alpha$ δεν έχει λύση όταν $\alpha\neq0$ και έχει άπειρες όταν $\alpha=0$.
\item Το $(-12):4$ είναι $-3$, όχι $3$: ετερόσημοι δίνουν αρνητικό.
\item Η διαίρεση δεν είναι αντιμεταθετική: $(-6):3=-2$ αλλά $3:(-6)=-\frac{1}{2}$.
\item Αντίστροφος δεν είναι αντίθετος: ο αντίθετος του $4$ είναι $-4$ και ο αντίστροφος $\frac{1}{4}$.
\end{itemize}
\end{prosoxi}

\tmima{Λυμένα παραδείγματα}

\paradeigma{Το ταμείο μοιράζεται}
\tixero Το ταμείο είναι $-24$ € και μοιράζεται σε $6$ παιδιά.\par
\tithelo Το μερίδιο του καθενός και πόσα παιδιά βάζουν από $-3$.\par
\possindeo $(-24):6=-4$ (ετερόσημοι, αρνητικό). $(-24):(-3)=+8$ (ομόσημοι, θετικό).\par
\apantisi $-4$ € και $8$ παιδιά.\par
\elegxos $6\cdot(-4)=-24$ και $8\cdot(-3)=-24$.

\paradeigma{Με τον ορισμό}
\tixero Το πηλίκο $12:(-3)$.\par
\tithelo Να το βρω ως λύση εξίσωσης.\par
\possindeo Ψάχνω $x$ με $(-3)\cdot x=12$. Αφού $(-3)\cdot(-4)=12$, είναι $x=-4$.\par
\apantisi $-4$.\par
\elegxos Ο κανόνας: ετερόσημοι, αρνητικό, $12:3=4$.

\paradeigma{Πηλίκα ρητών}
\tixero Τα $(-7{,}5):(-2{,}5)$, $\left(-\frac{3}{4}\right):\frac{1}{2}$, $(+0{,}6):(-1{,}5)$.\par
\tithelo Να τα υπολογίσω.\par
\possindeo Ομόσημοι: $7{,}5:2{,}5=3$, άρα $+3$. Ετερόσημοι: $-\frac{3}{4}\cdot2=-\frac{3}{2}$. Ετερόσημοι: $0{,}6:1{,}5=0{,}4$, άρα $-0{,}4$.\par
\apantisi $3$, $-\frac{3}{2}$, $-0{,}4$.\par
\elegxos $(-2{,}5)\cdot3=-7{,}5$ και $(-1{,}5)\cdot(-0{,}4)=0{,}6$.

\paradeigma{Αντίστροφοι}
\tixero Οι αριθμοί $-4$, $-\frac{2}{3}$, $5$, $0$.\par
\tithelo Οι αντίστροφοί τους.\par
\possindeo Ο αντίστροφος του $\alpha$ είναι ο $\frac{1}{\alpha}$ με το ίδιο πρόσημο. Του $-4$: $-\frac{1}{4}$. Του $-\frac{2}{3}$: $-\frac{3}{2}$. Του $5$: $\frac{1}{5}$. Το $0$ δεν έχει.\par
\apantisi $-\frac{1}{4}$, $-\frac{3}{2}$, $\frac{1}{5}$ και το $0$ δεν έχει αντίστροφο.\par
\elegxos $(-4)\cdot\left(-\frac{1}{4}\right)=1$.

\paradeigma{Παραστάσεις και σύνθετο κλάσμα}
\tixero Οι παραστάσεις $(-12):3+(-5)\cdot(-2)$ και $\dfrac{-\frac{6}{5}}{\frac{3}{10}}$.\par
\tithelo Να τις υπολογίσω.\par
\possindeo $(-12):3=-4$ και $(-5)(-2)=10$, άρα $-4+10=6$. Το σύνθετο: $-\frac{6}{5}\cdot\frac{10}{3}=-\frac{60}{15}=-4$.\par
\apantisi $6$ και $-4$.\par
\elegxos $-4\cdot\frac{3}{10}=-\frac{6}{5}$.

\tmima{Ασκήσεις}

\zoni{Το κατάλαβα}{πηλίκα, εξισώσεις, αντίστροφοι}

\ex Υπολόγισε: α) $(-24):6$ \quad β) $(-24):(-3)$ \quad γ) $18:(-6)$ \quad δ) $0:(-5)$ \quad ε) $(-7{,}5):(-2{,}5)$ \quad στ) $(+36):(+9)$.

\ex Λύσε: α) $(-3)x=12$ \quad β) $5x=-20$ \quad γ) $x:(-2)=6$ \quad δ) $(-4)x=-4$ \quad ε) $x:3=-5$ \quad στ) $12:x=-4$.

\ex Υπολόγισε: α) $\left(-\frac{3}{4}\right):\frac{1}{2}$ \quad β) $\left(-\frac{2}{3}\right):\left(-\frac{4}{9}\right)$ \quad γ) $5:\left(-\frac{2}{5}\right)$ \quad δ) $(-1{,}2):0{,}4$ \quad ε) $(+0{,}6):(-1{,}5)$ \quad στ) $\left(-\frac{1}{2}\right):\left(-\frac{1}{2}\right)$.

\ex Βρες τους αντίστροφους των $-4$, $-\frac{2}{3}$, $5$, $-0{,}5$, $\frac{1}{3}$.

\ex Απλοποίησε τα σύνθετα κλάσματα: α) $\dfrac{-\frac{6}{5}}{\frac{3}{10}}$ \quad β) $\dfrac{\frac{2}{3}}{-\frac{4}{9}}$ \quad γ) $\dfrac{-5}{-\frac{1}{2}}$.

\ex Υπολόγισε: α) $(-12):3+(-5)\cdot(-2)$ \quad β) $20:(-4)-(-3)\cdot2$ \quad γ) $[(-18):(-6)]\cdot(-4)$.

\ex Πες το πρόσημο χωρίς να υπολογίσεις: α) $(-45):9$ \quad β) $(-45):(-9)$ \quad γ) $45:(-9)$ \quad δ) $(-1\,000):(-1)$.

\ex Σωστό ή λάθος; α) $(-12):(-4)=-3$. \quad β) $0:(-5)=0$. \quad γ) $(-5):0=0$. \quad δ) $\alpha:\beta=\alpha\cdot\frac{1}{\beta}$. \quad ε) Το πηλίκο ετεροσήμων είναι αρνητικό. \quad στ) $\frac{-6}{3}=\frac{6}{-3}$.

\zoni{Το χρησιμοποιώ}{χρέη, ρυθμοί, μέσοι όροι}

\ex \textbf{Το ταμείο.} Το ταμείο είναι $-24$ € και μοιράζεται εξίσου στα $6$ παιδιά. Πόσο αντιστοιχεί στο καθένα; Αν κάθε παιδί βάζει $-3$ €, πόσα παιδιά έχουν συνεισφέρει συνολικά $-24$ €;

\ex \textbf{Η θερμοκρασία.} Η θερμοκρασία έπεσε $18$ βαθμούς σε $6$ ώρες. Ποιος είναι ο ρυθμός μεταβολής ανά ώρα; Και αν έπεσε $12$ βαθμούς σε $4$ ώρες;

\ex \textbf{Ο δύτης.} Ο δύτης κατέβηκε από το $0$ στο $-36$ μ. σε $9$ λεπτά. Πόσα μέτρα το λεπτό;

\ex \textbf{Σε πόση ώρα;} Η θερμοκρασία πέφτει $2{,}5$ βαθμούς την ώρα. Σε πόσες ώρες θα πέσει $15$ βαθμούς;

\ex \textbf{Ο μέσος όρος.} Οι θερμοκρασίες μιας εβδομάδας ήταν $-3$, $-5$, $2$, $0$, $-4$, $1$, $-1$. Βρες το άθροισμα και τον μέσο όρο στο δέκατο.

\ex \textbf{Οι πόντοι.} Οι διαφορές πόντων της Φανής σε $4$ αγώνες ήταν $-8$, $+4$, $-10$, $+6$. Ποιος είναι ο μέσος όρος;

\ex \textbf{Τα παιδιά.} Το χρέος των $-45$ € μοιράστηκε εξίσου και κάθε παιδί χρωστά $-5$ €. Πόσα είναι τα παιδιά;

\zoni{Το σκέφτομαι}{γιατί όχι με το μηδέν}

\ex Γιατί δεν διαιρούμε με το $0$; Εξέτασε με τον ορισμό τις $0\cdot x=5$ και $0\cdot x=0$.

\ex Ο Στέλιος γράφει $(-12):4=3$. Βρες το λάθος και έλεγξε το σωστό με πολλαπλασιασμό.

\ex Σύγκρινε τα πρόσημα των $(-6):3$, $6:(-3)$, $(-6):(-3)$.

\ex Αν $\alpha:\beta=2$ και $\alpha=-10$, βρες το $\beta$.

\ex Δείξε με $\alpha=6$ και $\beta=3$ ότι $\dfrac{-\alpha}{\beta}=\dfrac{\alpha}{-\beta}=-\dfrac{\alpha}{\beta}$.

\zoni{Για το σπίτι}{μοίρασε και σκέψου}

\ex Φτιάξε ένα πρόβλημα όπου ένα χρέος μοιράζεται ισόποσα και λύσε το.
